\DocumentMetadata{
  pdfversion=2.0,
  pdfstandard=ua-2,
  lang=en-US,
  tagging=on,
  tagging-setup={math/setup=mathml-SE,tabsorder=structure}
}
\documentclass[11pt,letterpaper]{article}
\usepackage[margin=0.68in,includefoot]{geometry}
\usepackage{newcomputermodern}
\usepackage{amsmath,amscd,mathtools}
\usepackage{tikz,adjustbox}
\providecommand{\square}{\mdlgwhtsquare}
\usepackage[hidelinks]{hyperref}
\hypersetup{
  pdftitle={Analysis PhD Exam, May 2022},
  pdfauthor={University of Florida Department of Mathematics},
  pdfsubject={Graduate mathematics examination},
  pdfdisplaydoctitle=true,
  bookmarksopen=true
}
\setcounter{secnumdepth}{0}
\setlength{\parindent}{0pt}
\setlength{\parskip}{0.28em}
\setlength{\emergencystretch}{3em}
\setlength{\leftmargini}{2.15em}
\setlength{\leftmarginii}{2.65em}
\setlength{\leftmarginiii}{3em}
\pagestyle{empty}
\raggedbottom
\allowdisplaybreaks
\NewTaggingSocketPlug{block/list/label}{examproblem}{%
  \tagstructbegin{tag=Lbl}\tagstructend%
  \tagstructbegin{tag=\LBody}%
  \tagstructbegin{tag=\ProblemHeadingTag,title-o={\ProblemHeadingText}}%
  \tagmcbegin{tag=\ProblemHeadingTag,actualtext={\ProblemHeadingText}}%
  #2%
  \tagmcend%
  \tagstructend%
}
\newenvironment{examproblems}[2]{%
  \def\ProblemHeadingTag{#2}%
  \AssignTaggingSocketPlug{block/list/label}{examproblem}%
  \begin{enumerate}%
  \setcounter{enumi}{#1}%
  \renewcommand{\labelenumi}{\textbf{\theenumi.}}%
  \setlength{\topsep}{0.35em}%
  \setlength{\partopsep}{0pt}%
  \setlength{\itemsep}{0.48em}%
  \setlength{\parsep}{0.18em}%
}{\end{enumerate}}
\newenvironment{examsubparts}{%
  \AssignTaggingSocketPlug{block/list/label}{default}%
  \begin{enumerate}%
  \setlength{\topsep}{0.18em}%
  \setlength{\partopsep}{0pt}%
  \setlength{\itemsep}{0.16em}%
  \setlength{\parsep}{0.1em}%
}{\end{enumerate}}
\newenvironment{examinstructions}{%
  \par\begingroup\itshape
}{%
  \par\endgroup\vspace{0.18em}
}
\makeatletter
\renewcommand\subsection{\@startsection{subsection}{2}{0pt}%
  {-1.25ex plus -.25ex minus -.1ex}%
  {0.55ex plus .1ex}%
  {\normalfont\large\bfseries}}
\renewcommand{\@maketitle}{%
  \newpage\null\vskip -1em%
  {\footnotesize\bfseries
    \href{https://www.ufl.edu/}{University of Florida}, \href{https://math.ufl.edu/}{Department of Mathematics}\par}%
  \vskip -0.5em%
  \tagstructbegin{tag=H1,title={Analysis PhD Exam, May 2022}}%
  \tagpdfparaOff%
  \tagmcbegin{tag=H1}%
    {\LARGE\bfseries\href{https://gma.math.ufl.edu/exams/analysis-phd/}{Analysis PhD Exam}, May 2022\par}%
  \tagmcend%
  \tagpdfparaOn%
  \tagstructend%
  \vskip 0.25em%
  \tagmcbegin{artifact}%
    \hrule height 1pt%
  \tagmcend%
  \vskip 1em%
}
\makeatother
\title{Analysis PhD Exam, May 2022}
\author{}
\date{}
\begin{document}
\maketitle
\thispagestyle{empty}
\begin{examinstructions}
Write solutions in a neat and logical fashion, giving complete reasons for all steps. Attempt SIX problems.
\end{examinstructions}
\begin{examproblems}{0}{H2}
\def\ProblemHeadingText{Problem 1.}
\item
State Tonelli’s theorem, and give an example to show that the~\(\sigma\)-finiteness hypothesis is necessary.
\def\ProblemHeadingText{Problem 2.}
\item
Let~\(\mu\) be a positive measure and let \((E_n)\) be a sequence of measurable sets with \(\mu(E_n)<\infty\) for all~\(n\). Prove that if \(\mathbf{1}_{E_n}\to f\) in \(L^1(\mu)\), then there exists a measurable set~\(E\) so that \(f=\mathbf{1}_E\) a.e. and \(\mu(E)=\lim \mu(E_n)\).
\def\ProblemHeadingText{Problem 3.}
\item
Let~\(\mu\) and~\(\nu\) be finite, positive measures defined on the same measurable space \((X,\mathcal{M})\), and suppose that \(\nu\ll\mu\).
\begin{examsubparts}
\item[\textnormal{a)}]
Prove that if \(L^1(\mu)\subset L^1(\nu)\), then the inclusion map \(\iota:L^1(\mu)\hookrightarrow L^1(\nu)\) is necessarily bounded.
\item[\textnormal{b)}]
Prove that if \(L^1(\mu)\subset L^1(\nu)\) then \(\frac{d\nu}{d\mu}\in L^\infty(\mu)\).
\end{examsubparts}
\def\ProblemHeadingText{Problem 4.}
\item
Let~\(\mu\) be a positive measure and suppose that \(1<p<r<q<\infty\). Prove that
\[
L^r(\mu)\subset L^p(\mu)+L^q(\mu).
\]
 (That is, every \(f\in L^r(\mu)\) can be written as a sum \(f=g+h\) where \(g\in L^p(\mu)\) and \(h\in L^q(\mu)\).)
\def\ProblemHeadingText{Problem 5.}
\item
Let~\(\mathcal{X}\) be a Banach space, and let \(\mathcal{M},\mathcal{N}\) be closed subspaces of~\(\mathcal{X}\). Suppose that every \(x\in\mathcal{X}\) can be decomposed uniquely as
\[
x=m+n
\]
 with \(m\in\mathcal{M}\), \(n\in\mathcal{N}\). Prove that the assignment \(x\to m\) defines a bounded linear operator from~\(\mathcal{X}\) to~\(\mathcal{M}\).
\def\ProblemHeadingText{Problem 6.}
\item
This problem has two parts.
\begin{examsubparts}
\item[\textnormal{a)}]
Prove that for each integer \(n\ge 0\), there exists a function \(f_n\in L^2[0,1]\) such that for every polynomial~\(p\) of degree at most~\(n\),
\[
p(1)=\int_0^1 p(x)f_n(x)\,dx.
\]
 Is \(f_n\) unique?
\item[\textnormal{b)}]
Is there a function \(f\in L^2[0,1]\) such that
\[
p(1)=\int_0^1 p(x)f(x)\,dx
\]
 for all polynomials~\(p\)? Prove or disprove.
\end{examsubparts}
\def\ProblemHeadingText{Problem 7.}
\item
This problem has two parts.
\begin{examsubparts}
\item[\textnormal{a)}]
State the Banach Isomorphism Theorem and the Closed Graph Theorem.
\item[\textnormal{b)}]
Using the Banach Isomorphism Theorem (or otherwise), prove the Closed Graph Theorem.
\end{examsubparts}
\def\ProblemHeadingText{Problem 8.}
\item
(Short answer, you do not need to give a detailed proof, just a brief explanation.)
\begin{examsubparts}
\item[\textnormal{a)}]
Define the Fourier transform on \(L^1(\mathbb{R})\) and sketch its extension to \(L^2(\mathbb{R})\).
\item[\textnormal{b)}]
Sketch a construction of a bounded linear functional~\(\lambda\) on \(L^\infty(\mathbb{R})\) which is not of the form \(\lambda(f)=\int f(x)g(x)\,dx\) for any \(L^1\) function~\(g\).
\item[\textnormal{c)}]
Sketch a proof of the fact that there is no norm under which \(c_{00}\) is complete.
\end{examsubparts}
\end{examproblems}
\end{document}
